% ---------------- 第5章 章末 ----------------
\zhangmo{本章小结}
\subsection*{一、知识总结}
\begin{center}
\small
\begin{tabularx}{\linewidth}{p{1.6cm}XX}
\toprule
& 原电池 & 电解池 \\
\midrule
能量转化 & 化学能$\to$电能（自发） & 电能$\to$化学能（非自发也能进行） \\
电极名称 & 负极（氧化）、正极（还原） & 阳极（氧化，接电源正极）、阴极（还原，接电源负极） \\
离子迁移 & 阳离子$\to$正极，阴离子$\to$负极 & 阳离子$\to$阴极，阴离子$\to$阳极 \\
典型应用 & 化学电源、加快反应、牺牲阳极法 & 氯碱、电镀、精炼、电冶金、外加电流法 \\
\bottomrule
\end{tabularx}
\end{center}

\subsection*{二、方法提升}
\begin{enumerate}
\item \textbf{判电极}：看化合价变化（升高为负极/阳极），或看外接电源，或看现象（溶解、析出、放气）。
\item \textbf{写反应}：先总后分、按介质补粒子；燃料电池“负极$=$总反应$-$正极”。
\item \textbf{判膜}：从电荷平衡和产品需求出发，确定哪种离子必须通过。
\item \textbf{算数量}：抓住电子守恒，建立 $e^-$ 与各物质的关系式。
\end{enumerate}

\subsection*{三、易错点归纳}
\begin{enumerate}
\item 认为活泼金属一定作负极（Mg–Al–NaOH、Fe–Cu–浓硝酸是反例）。
\item 碱性介质中电极反应出现 \ce{H+} 或 \ce{CO2}；酸性介质中出现 \ce{OH-}。
\item 认为电子可以通过电解质溶液；把离子迁移方向弄反。
\item 电解时忽略活性阳极自身溶解；电解后复原“加错物质”。
\item 外加电流法中把被保护金属接电源正极（应接负极）。
\end{enumerate}

\zhangmo{分层练习}
\zu{A组\quad 基础巩固}
\begin{xiti}
\item 关于 Zn–Cu 双液原电池（盐桥中为 KCl 琼胶），下列说法正确的是\xz
\xxa{电子由锌片经盐桥流向铜片}{盐桥中的 \ce{K+} 移向 \ce{ZnSO4} 溶液}{铜片上发生的反应为 \ce{Cu^2+ + 2e- = Cu}}{取出盐桥，电流计指针仍发生偏转}

\item 将镁片和铝片用导线相连插入 NaOH 溶液中，下列说法正确的是\xz
\xxa{镁比铝活泼，镁作负极}{负极反应为 \ce{Al - 3e- + 4OH- = [Al(OH)4]-}}{正极上产生 \ce{O2}}{电子由镁片经导线流向铝片}

\item 以 KOH 溶液为电解质的甲烷燃料电池中，负极反应式为\xz
\xxa{\ce{CH4 - 8e- + 2H2O = CO2 + 8H+}}{\ce{CH4 - 8e- + 10OH- = CO3^2- + 7H2O}}{\ce{O2 + 4e- + 2H2O = 4OH-}}{\ce{CH4 + 8e- + 10OH- = CO3^2- + 7H2O}}

\item 关于铅酸蓄电池，下列说法正确的是\xz
\xxa{放电时，负极质量减小}{放电时，电解质溶液的密度增大}{充电时，铅酸蓄电池的负极应与外接电源的正极相连}{充电时，阳极反应为 \ce{PbSO4 - 2e- + 2H2O = PbO2 + 4H+ + SO4^2-}}

\item 用惰性电极电解足量 \ce{CuSO4} 溶液一段时间后，要使溶液恢复到原来的浓度和体积，应加入适量的\xz
\xx{Cu}{CuO}{\ce{Cu(OH)2}}{\ce{CuSO4}}

\item \laiyuan{2024年广东卷第5题，来源：仓库所附真题文档}\ 我国自主设计建造的浮式生产储卸油装置“海葵一号”将在珠江口盆地海域使用，其钢铁外壳镶嵌了锌块，以利用电化学原理延缓外壳的腐蚀。下列有关说法正确的是\xz
\xxb{钢铁外壳为负极}{镶嵌的锌块可永久使用}{该法为外加电流法}{锌发生反应：\ce{Zn - 2e- = Zn^2+}}

\item 用铜作电极电解稀硫酸，下列说法正确的是\xz
\xxa{阳极上放出 \ce{O2}}{阴极上析出 Cu}{溶液逐渐变蓝，阴极上放出 \ce{H2}}{电解一段时间后溶液 pH 减小}
\end{xiti}

\zu{B组\quad 高考典型}
\begin{xiti}
\item \laiyuan{2025年广东卷第16题，来源：仓库所附真题文档}\ 某理论研究认为：氢氧燃料电池（图b，两电极间为质子交换膜）的电极Ⅰ和Ⅱ上所发生反应的催化机理分别如图a和图c。图a：\ce{O2} 在电极Ⅰ表面的催化剂上依次经过 i$\to$ii$\to$iii$\to$iv 四步，每步结合 1 个 \ce{H+} 和 1 个 $e^-$，最终生成 \ce{H2O}；图c：\ce{H2} 在电极Ⅱ表面经两步各失去 1 个电子生成 \ce{H+}。其中 \ce{O2} 获得第一个电子的过程（i$\to$ii）最慢。由此可知，理论上\xz
\xxb{负极反应的催化剂是 i}{图a中，i 到 ii 过程的活化能一定最低}{电池工作过程中，负极室的溶液质量保持不变}{相同时间内，电极Ⅰ和电极Ⅱ上的催化循环完成次数相同}

\item \laiyuan{原创}\ 磷酸铁锂电池的总反应为 \ce{Li_xC6 + Li_{1-$x$}FePO4 <=>[\text{放电}][\text{充电}] 6C + LiFePO4}，电解质为能传导 \ce{Li+} 的有机溶液。(1) 写出放电时正极的电极反应式。(2) 充电时，\ce{Li+} 向哪一极迁移？阳极反应式是什么？(3) 放电时每转移 1 mol 电子，负极材料质量变化多少？(4) 为什么该电池不能用水溶液作电解质？

\item 下列关于电镀和电解精炼的说法正确的是\xz
\xxa{在铁件上镀锌，铁件作阳极，锌片作阴极}{电镀过程中，电解质溶液中镀层金属离子的浓度明显减小}{电解精炼铜时，粗铜作阳极，阳极减少的质量等于阴极增加的质量}{电解精炼铜时，Ag、Au 等不如铜活泼的金属沉积在阳极底部形成阳极泥}

\item \laiyuan{原创}\ 用如下三室装置电解 \ce{Na2SO4} 溶液制备 NaOH 和 \ce{H2SO4}：左侧为阳极室（稀硫酸），中间室为 \ce{Na2SO4} 溶液，右侧为阴极室（稀 NaOH 溶液），两张膜分别为 X 和 Y（X 在阳极室与中间室之间）。(1) X、Y 分别是阴离子交换膜还是阳离子交换膜？(2) 写出阳极、阴极的电极反应式。(3) 以 2.0 A 电流电解 965 s，理论上阴极室生成 NaOH 的物质的量为多少？中间室 \ce{Na2SO4} 减少多少克？

\item \laiyuan{原创}\ 锌—空气电池以 KOH 溶液为电解质，总反应为 \ce{2Zn + O2 + 4OH- + 2H2O = 2[Zn(OH)4]^2-}。(1) 写出正、负极的电极反应式。(2) 放电时，\ce{OH-} 向哪一极移动？电解质溶液中 $c(\ce{OH-})$ 如何变化？(3) 每消耗 6.5 g Zn，理论上消耗标准状况下空气多少升（空气中 \ce{O2} 体积分数按 20\% 计）？

\item \laiyuan{原创}\ 将两个电解池串联：甲池为惰性电极电解足量 \ce{AgNO3} 溶液，乙池为电解精炼铜（粗铜含少量 Zn、Ag）。通电一段时间后甲池阴极增重 10.8 g。(1) 乙池阴极增重多少克？(2) 乙池阳极减少的质量是否等于阴极增加的质量？为什么？(3) 甲池阳极产生气体在标准状况下的体积为多少？
\end{xiti}

\zu{C组\quad 冲刺高分}
\begin{xiti}
\item \laiyuan{2026年广东卷第16题，来源：仓库所附真题文档（网络收集版）}\ 一种利用电解还原 \ce{CO2} 同时实现海水脱盐的装置：从左到右依次为电极a、Ⅰ室、膜①、Ⅱ室（海水）、膜②、Ⅲ室（海水）、膜③、Ⅳ室、电极b；Ⅰ室中通入 \ce{CO2} 并放出 CO，Ⅳ室放出 \ce{O2}；电解液在Ⅰ和Ⅳ室间循环，初始电解液为 \ce{NaHCO3} 溶液；膜①～③为阳离子或阴离子交换膜。下列说法错误的是\xz
\xxb{①为阳离子交换膜，②为阴离子交换膜}{阳极反应为 \ce{4OH- - 4e- = 2H2O + O2 ^}}{电解一段时间后，Ⅲ室海水中盐的浓度降低}{电解过程中，Ⅰ室电解液的 pH 在一段时间内可保持稳定}

\item \laiyuan{原创}\ 以质子交换膜甲醇燃料电池为电源，用惰性电极电解 200 mL 0.50 mol·L$^{-1}$ \ce{CuSO4} 溶液。(1) 写出燃料电池负极、正极的电极反应式。(2) 当燃料电池消耗 0.32 g 甲醇时，电解池阴极析出铜多少克？阳极产生气体在标准状况下的体积为多少？(3) 此时电解液中 $c(\ce{H+})$ 为多少（忽略体积变化）？要使电解液复原，应加入什么物质、多少克？(4) 若继续电解至阴极开始产生气体，此后电解的实质是什么？

\item \laiyuan{原创}\ 用电解法将 \ce{Na2CrO4} 转化为 \ce{Na2Cr2O7}：电解槽用阳离子交换膜隔开，阳极室加入 \ce{Na2CrO4} 溶液，阴极室加入稀 NaOH 溶液，两极均为惰性电极。已知 \ce{2CrO4^2- + 2H+ <=> Cr2O7^2- + H2O}。(1) 写出阳极、阴极的电极反应式。(2) 解释阳极室中 \ce{Na2CrO4} 转化为 \ce{Na2Cr2O7} 的原理。(3) 通过阳离子交换膜迁移的离子是什么？向哪个方向迁移？(4) 若阳极室溶液中 Na 与 Cr 的物质的量之比由开始的 2:1 变为 $d$:1，则此时 \ce{Na2CrO4} 的转化率为\tk（用 $d$ 表示）。
\end{xiti}

\zhangmo{练习解析}
\begin{jiexilist}
\item \ans{C}电子只经导线流动，A 错；盐桥中阳离子 \ce{K+} 移向正极区（\ce{CuSO4} 溶液），B 错；铜为正极，\ce{Cu^2+} 得电子，C 对；取出盐桥，回路断开，无电流，D 错。
\item \ans{B}Mg 不与 NaOH 反应，Al 与 NaOH 反应，Al 作负极，A 错、B 对；正极上 \ce{H2O} 得电子生成 \ce{H2}，C 错；电子由负极（铝）流向正极（镁），D 错。
\item \ans{B}碱性介质中 C 元素以 \ce{CO3^2-} 形式存在，不能写 \ce{H+} 和 \ce{CO2}；C 是正极反应；D 中负极应失电子。
\item \ans{D}放电时负极 Pb 生成 \ce{PbSO4}，质量增大，A 错；消耗 \ce{H2SO4}、生成水，溶液密度减小，B 错；充电时“负接负”，C 错；充电阳极是放电正极反应的逆反应，D 对。
\item \ans{B}电解时阴极析出 Cu，阳极放出 \ce{O2}，“出什么加什么”，相当于减少了 CuO，应加 CuO；加 \ce{Cu(OH)2} 会多加水。
\item \ans{D}Zn 比 Fe 活泼，锌块为负极，钢铁外壳为正极被保护，A 错；锌不断溶解，需定期更换，B 错；该法为牺牲阳极法，C 错；D 对。
\item \ans{C}阳极 Cu 为活性电极，自身溶解：\ce{Cu - 2e- = Cu^2+}，溶液变蓝；阴极开始时 \ce{H+} 放电产生 \ce{H2}（\ce{H+} 浓度大，生成的 \ce{Cu^2+} 较少），总反应 \ce{Cu + H2SO4 ->[\text{电解}] CuSO4 + H2 ^}，消耗 \ce{H+}，pH 增大。
\item \ans{C}电极Ⅰ上 \ce{O2} 被还原，为正极，i 是正极反应的催化剂，A 错；i$\to$ii 最慢，活化能最高，B 错；负极 \ce{H2 - 2e- = 2H+}，每失去 1 个电子就有 1 个 \ce{H+} 通过质子交换膜进入正极室，负极室溶液质量不变，C 对；\ce{O2} 完成一次循环需 4 个电子，\ce{H2} 完成一次循环需 2 个电子，同一时间转移电子数相同，两者循环次数之比为 1:2，D 错。
\item (1) \ce{Li_{1-$x$}FePO4 + $x$Li+ + $x$e- = LiFePO4}。(2) 充电时阳离子移向阴极（即放电时的负极，石墨一侧）；阳极：\ce{LiFePO4 - $x$e- = Li_{1-$x$}FePO4 + $x$Li+}。(3) 负极 \ce{Li_xC6 - $x$e- = $x$Li+ + 6C}，每转移 1 mol 电子，负极减少 1 mol Li，即质量减少 7 g。(4) Li 及嵌锂石墨很活泼，会与水剧烈反应。
\item \ans{D}镀锌时铁件作阴极、锌片作阳极，A 错；电镀时阳极溶解的金属离子与阴极析出的相等，浓度基本不变，B 错；精炼时阳极上 Zn、Fe、Ni 等也溶解且有阳极泥，两极质量变化不相等，C 错；D 对。
\item (1) 阳极室需要 \ce{SO4^2-} 进入（与生成的 \ce{H+} 组成 \ce{H2SO4}），X 为阴离子交换膜；阴极室需要 \ce{Na+} 进入（与生成的 \ce{OH-} 组成 NaOH），Y 为阳离子交换膜。(2) 阳极：\ce{2H2O - 4e- = O2 ^ + 4H+}；阴极：\ce{2H2O + 2e- = H2 ^ + 2OH-}。(3) $n(e^-)=\dfrac{2.0\times965}{96500}=0.020$ mol；生成 NaOH 0.020 mol；中间室迁出 0.020 mol \ce{Na+} 和 0.010 mol \ce{SO4^2-}，即减少 0.010 mol \ce{Na2SO4}，1.42 g。
\item (1) 负极：\ce{Zn - 2e- + 4OH- = [Zn(OH)4]^2-}；正极：\ce{O2 + 4e- + 2H2O = 4OH-}。(2) 移向负极；总反应消耗 \ce{OH-}，$c(\ce{OH-})$ 减小。(3) 0.1 mol Zn 失 0.2 mol 电子，需 0.05 mol \ce{O2}，即 1.12 L，空气 $1.12\div20\%=5.6$ L。
\item (1) 甲池析出 Ag 0.1 mol，转移 0.1 mol 电子；乙池阴极析出 Cu 0.05 mol，增重 3.2 g。(2) 不相等。粗铜中的 Zn 也失电子溶解（摩尔质量与 Cu 不同），Ag 不溶解而落入阳极泥，所以阳极减少的质量与阴极增加的质量不相等，但两极转移的电子数相等。(3) 甲池阳极：\ce{2H2O - 4e- = O2 ^ + 4H+}，0.1 mol 电子生成 0.025 mol \ce{O2}，即 0.56 L。
\item \ans{C}a 极 \ce{CO2} 得电子生成 CO，为阴极：\ce{CO2 + H2O + 2e- = CO + 2OH-}，生成的 \ce{OH-} 与 \ce{HCO3-} 反应生成 \ce{CO3^2-}；b 极为阳极，水电离的 \ce{OH-} 放电生成 \ce{O2}，B 对。若膜③为阴离子交换膜，Ⅲ室 \ce{Cl-} 会进入Ⅳ室在阳极放电生成 \ce{Cl2}，故③应为阳离子交换膜（Ⅳ室 \ce{Na+} 进入Ⅲ室）；要使某室脱盐，该室的 \ce{Na+}、\ce{Cl-} 必须都能移出：Ⅱ室 \ce{Na+} 经膜①（阳膜）移入Ⅰ室，\ce{Cl-} 经膜②（阴膜）移入Ⅲ室，A 对。于是Ⅱ室盐浓度降低（脱盐），Ⅲ室得到 \ce{Cl-}（来自Ⅱ）和 \ce{Na+}（来自Ⅳ），盐浓度升高，C 错。Ⅰ室中 \ce{HCO3-}/\ce{CO3^2-} 构成缓冲体系，且与Ⅳ室电解液循环，pH 可在一段时间内保持稳定，D 对。
\item (1) 负极：\ce{CH3OH - 6e- + H2O = CO2 ^ + 6H+}；正极：\ce{O2 + 4e- + 4H+ = 2H2O}。(2) $n(\ce{CH3OH})=0.01$ mol，转移 0.06 mol 电子。$n(\ce{Cu^2+})=0.10$ mol，足量，析出 Cu 0.03 mol，即 1.92 g；阳极 \ce{O2} 0.015 mol，0.336 L。(3) 生成 \ce{H+} 0.06 mol，$c(\ce{H+})=0.30$ mol·L$^{-1}$；溶液中减少的是 0.03 mol Cu 和 0.015 mol \ce{O2}，相当于 0.03 mol CuO，加入 CuO 2.4 g（或 \ce{CuCO3} 3.72 g）。(4) \ce{Cu^2+} 完全放电后，阴极上 \ce{H+}（来自水和生成的 \ce{H2SO4}）放电，阳极仍放出 \ce{O2}，实质是电解水。
\item (1) 阳极：\ce{2H2O - 4e- = O2 ^ + 4H+}；阴极：\ce{2H2O + 2e- = H2 ^ + 2OH-}。(2) 阳极生成 \ce{H+}，使 \ce{2CrO4^2- + 2H+ <=> Cr2O7^2- + H2O} 正向移动。(3) \ce{Na+}，由阳极室迁移到阴极室（阴极室得到 NaOH）。(4) $2-d$。设起始 \ce{Na2CrO4} 1 mol（Na 2 mol，Cr 1 mol），转化率为 $\alpha$：生成 $\alpha$ mol \ce{H+} 需转移 $\alpha$ mol 电子，同时有 $\alpha$ mol \ce{Na+} 迁出阳极室，剩余 Na $(2-\alpha)$ mol，Cr 仍为 1 mol。$\dfrac{2-\alpha}{1}=d$，$\alpha=2-d$。
\end{jiexilist}
